\documentclass[11pt]{amsart}
\usepackage[margin=1.05in]{geometry}
\usepackage{amsmath,amssymb}
\usepackage[colorlinks=true,linkcolor=blue,citecolor=blue,urlcolor=blue]{hyperref}

\title{The \(1\le p\le2\) Translate-Frame Gap After Freeman--Odell--Schlumprecht--Zsak}
\author{fa\_banach\_001 literature-status packet}
\date{2026-07-03}

\begin{document}
\maketitle

\section*{Status}

This is a literature-already-answered packet.  The source paper is
Freeman--Odell--Schlumprecht--Zsak, \emph{Unconditional structures of
translates for \(L_p(\mathbb R^d)\)}, arXiv:1209.4619.  In Section 3,
Theorem 3.2, the authors prove that if \(2<p<\infty\), \(d\in\mathbb N\),
and \((\lambda_n)\) is any unbounded sequence in \(\mathbb R^d\), then there
is an unconditional Schauder frame
\[
  (T_{\lambda_n}f,g_n^*)_{n\in\mathbb N}
\]
for \(L_p(\mathbb R^d)\).  This leaves the complementary arbitrary-translate
frame case \(1<p\le2\) as the missing half of the dichotomy.

\section*{Answer Source}

Lev--Tselishchev, \emph{There are no unconditional Schauder frames of
translates in \(L^p(\mathbb R)\), \(1\le p\le2\)}, arXiv:2312.01757,
explicitly identify this gap.  Their introduction states that the existence
of an unconditional Schauder frame of translates in \(L^p(\mathbb R)\),
\(1<p\le2\), remained open after the Freeman--Odell--Schlumprecht--Zsak
positive result for \(p>2\).  Their Theorem 1.1 proves that, for
\(1\le p\le2\), there is no unconditional Schauder frame of the form
\[
  (T_{\lambda_n}g,g_n^*)_{n\in\mathbb N}
\]
in \(L^p(\mathbb R)\).

\section*{Identification}

Combining the two papers gives the sharp line-case dichotomy for
unconditional Schauder frames whose vector side consists of translates of a
single function:
\[
  \text{existence exactly for } 2<p<\infty.
\]
The supporting authors explicitly know they are addressing the remaining
\(1<p\le2\) case after arXiv:1209.4619, so this is an
already-known literature answer, not a new proof from this run.

\section*{Scope}

This does not settle the long-standing Schauder-basis problem for translates
in \(L_p(\mathbb R)\).  It also does not settle the final Section 6 questions
in arXiv:1209.4619 about conditional bases/basic sequences or blocking into
unconditional FDDs.  The resolved item is the unconditional Schauder
translate-frame gap for \(1\le p\le2\).

\section*{References}

\begin{thebibliography}{9}
\bibitem{FOSZ}
D. Freeman, E. Odell, Th. Schlumprecht, and A. Zs\'ak,
\emph{Unconditional structures of translates for \(L_p(\mathbb R^d)\)},
arXiv:1209.4619.

\bibitem{LT}
N. Lev and A. Tselishchev,
\emph{There are no unconditional Schauder frames of translates in
\(L^p(\mathbb R)\), \(1\le p\le2\)},
arXiv:2312.01757.
\end{thebibliography}

\end{document}
